% BidPilot (GigBid) - R&D report. Compile: latexmk -pdf report.tex  (or upload the report/ folder to Overleaf).
% Every number is generated from evals/results/*.json by `python -m scripts.build_report_tables` into report/generated/.
\documentclass[10pt,a4paper]{article}
\usepackage[margin=2cm]{geometry}
\usepackage[T1]{fontenc}
\usepackage{lmodern,microtype,booktabs,xcolor,enumitem,graphicx,caption,float}
\usepackage[hidelinks]{hyperref}
\setlist{nosep,leftmargin=*}
\captionsetup{font=small,labelfont=bf}
\newcommand{\pending}{\textcolor{orange!80!black}{\textit{pending}}}
\newcommand{\pendingtable}[2]{\begin{center}\fbox{\parbox{0.9\linewidth}{\small\textbf{#1:} measurement \pending{} -- regenerate with
  \texttt{make eval \&\& python -m scripts.build\_report\_tables}; source file \texttt{#2}.}}\end{center}}
\newcommand{\file}[1]{\texttt{\detokenize{#1}}}
\InputIfFileExists{generated/numbers.tex}{}{}
\providecommand{\ragvzeroRfive}{\pending}\providecommand{\ragbestRfive}{\pending}\providecommand{\ragvtwoRfive}{\pending}
\providecommand{\ragvthreeRfive}{\pending}\providecommand{\ragveightRfive}{\pending}
\providecommand{\decF}{\pending}\providecommand{\decScamP}{\pending}\providecommand{\decScamR}{\pending}
\providecommand{\decInj}{\pending}\providecommand{\decPrice}{\pending}\providecommand{\decPninetyfive}{\pending}

\title{\textbf{BidPilot: a multi-agent freelancer copilot}\\[2pt]\large R\&D report (AI-4016, GigBid) -- engineering design review}
\author{Qasim Hafeez (22F-3812)\\\small Repository: \url{https://github.com/qasim-hafeez00/bidpilot}}
\date{November 2026}

\begin{document}
\maketitle
\begin{abstract}\noindent
BidPilot is the GigBid system: LangGraph agents find freelance jobs, decide whether they are worth a bid, write proposals
grounded in portfolio evidence and handle client conversations, while a human approves every proposal, message and
Connects spend. All marketplace access goes through one port with a GigSim adapter and a live Upwork MCP adapter. The
system runs end to end on a CPU laptop with open-weight models (Ollama, Qwen2.5 3B/7B) and CPU-friendly ONNX retrieval
models. This report states what was built, what was measured, and, just as importantly, what was not measured
and why. Every number below is generated from a file in \file{evals/results/}; cells marked \pending{} were not
measured at the commit this PDF was built from.
\end{abstract}

\section{Problem and constraints}
\textbf{Workflow.} A freelancer scans feeds, spends paid Connects, writes proposals and handles replies. BidPilot automates
discovery (Scout), assessment (Analyst, Risk, Fit, Pricing) and drafting (Writer, Critic, Compliance, Liaison). The
freelancer keeps every decision that leaves the system.

\textbf{Never automated (enforced in code, not in prompts).} Submitting a proposal, sending a message and spending
Connects all pass through a LangGraph \texttt{interrupt()} (\file{agent/nodes/approval.py},
\file{agent/nodes/client_liaison.py}). Per-agent MCP allowlists mean no agent except the post-approval
\texttt{submitter} can even see \texttt{submit\_proposal} (\file{agent/mcp_client.py}). The optional autonomy setting can only
auto-\emph{skip} low-fit jobs. Live mode only creates drafts that the user confirms on Upwork.

\textbf{Budgets (declared in \file{agent/config.yaml}).} p95 latency per job analysed $\le$ 20\,s and per proposal $\le$ 300\,s
on the open (CPU) profile. Cost per job analysed $\le$ \$0.002 and per proposal $\le$ \$0.02 on the hosted-open profile; the
local profile costs \$0 in API fees. Graph budgets: 30 steps, 16k tokens and at most 2 Critic revisions per thread.

\textbf{Compute reality.} All development and evaluation ran on an Intel i7-8550U (4 cores, 16\,GB RAM, no GPU). Measured
Ollama throughput was 6.1 tok/s generation for Qwen2.5-3B and 2.9 tok/s for Qwen2.5-7B (\S\ref{sec:models}). This one
constraint shaped the design: rules-first nodes, LLMs only where they add value, an LLM cache, and reduced sample
sizes for LLM-heavy experiments. Every reduction is stated where it occurs.

\section{Data}
The instructor's GigSim, persona corpora and dev set were not part of the starter kit available to me, so I rebuilt
them (\file{scripts/build_data/}), keeping real and synthetic data distinct. Every record carries a \texttt{provenance} field;
full details are in \file{docs/DATA_CARD.md}.
\begin{itemize}
\item \textbf{Real:} 27{,}713 usable Upwork postings with Upwork skill tags, client statistics (payment verification, spend,
  rating, reviews, hires, country) and Connects cost (Kaggle \emph{upwork-freelance-jobs-60k}). Pricing statistics come
  from the 53k-posting \emph{upwork-job-postings-dataset-2024} (percentiles per domain and rate type). Scam phrasing comes
  from the labelled fraudulent postings in \emph{trial-1}.
\item \textbf{Synthetic:} three personas (AI/ML, full-stack, data analyst), each with a PDF CV, 17--18 case studies, repositories
  (README plus code), reviews and 150 past proposals labelled won/replied/ignored. The corpus is built to defeat naive RAG:
  each project has \emph{implied} skills that appear only in code imports (e.g.\ \texttt{from celery import shared\_task}) and
  are paraphrased in prose (``a background worker queue''), and each project is described three ways (case study, README,
  proposal). Client hidden preferences, scam payloads, injected instructions and screening questions are also synthetic.
\item \textbf{Rejected:} the \emph{freelancer-success} dataset, because its hourly rates are capped (median \$280 for every
  skill), and the 1.3M-row contracts dataset, which is redundant with the postings for pricing.
\item \textbf{Splits:} a dev set of 200 jobs per persona with gold bid/skip labels, scam and injection flags, graded gold
  evidence and a fair price band. A held-out split of 80 jobs per persona with new clients is git-ignored and stands in
  for the instructor's hidden set.
\end{itemize}
\textbf{Threat to validity.} I wrote both the persona generator and the system, so the corpus shares my vocabulary.
Results here are therefore an upper bound on what a stranger's portfolio would yield.

\section{Architecture}
\begin{figure}[H]\centering\small
\begin{verbatim}
 Scout (worker, cron) --job_id--> intake --Send--> Job Analyst ----\
                                         --Send--> Risk Assessor ---> decide --skip/watch--> END
                                         --Send--> Fit Matcher -----/     \--bid--> Pricing -> Writer -> verify claims
   Critic --fail (<=2x)--> Writer;  Critic --pass--> Compliance --block--> END (reason shown in UI)
   Compliance --> dry-run preview --> [interrupt: approve | edit | reject] --> submit (idempotent) --> END
 Conversation graph: read messages -> classify intent -> (slots) -> draft reply -> [interrupt] -> send (+ tentative hold)
\end{verbatim}
\caption{Both graphs are hand-written \texttt{StateGraph}s (\file{agent/graphs/}); rendered Mermaid sources are in \file{docs/images/}.}
\end{figure}
\textbf{Why each agent exists.} The Analyst produces a structured \texttt{JobSpec} (rules first, optional small LLM). The Risk
Assessor needs the client history tool and a lexicon calibrated on scam text. The Fit Matcher is the only reader of the
portfolio, and it maps each decomposed requirement to evidence. Pricing combines market percentiles, the client range and
the user's floor with a Connects expected-value check. The Writer is the only node that needs the strongest model. The
Critic and Compliance Guard are separate because they fail for different reasons (quality vs.\ policy) and only
Compliance may block. The Liaison owns conversations.

\textbf{State, trace and memory are kept apart} (\file{agent/state.py}). \emph{State} holds typed fields with reducers wherever the
parallel branches write (\texttt{analyses}, \texttt{usage}, \texttt{trace}, and a severity-ordered \texttt{status} reducer). The
\emph{execution trace} is append-only \texttt{trace} records plus LangSmith/Phoenix spans plus the MCP audit log, with one trace id
from the UI request to the audit row. \emph{Long-term memory} is a LangGraph store with two namespaces, written only through
\file{agent/memory/promotion_rules.py}. A style rule is promoted when an edit pattern (e.g.\ \texttt{removed\_word:passionate},
\texttt{cut\_first\_sentence}) appears in $\ge$3 of the user's last 5 edits, and demoted when it appears in none. An outcome
insight is promoted only when both opening-style arms have $n\ge5$ and their reply rates differ by $\ge$15 points.

\textbf{Hand-offs} are Pydantic models (\file{agent/handoffs.py}): \texttt{JobSpec}, \texttt{RiskReport}, \texttt{FitReport}
(with a requirement-to-evidence matrix), \texttt{BidDecision}, \texttt{PriceQuote}, \texttt{ProposalDraft}, \texttt{Critique},
\texttt{ComplianceVerdict}, \texttt{Approval}.

\textbf{MCP surface} (\file{docs/TOOLS.md}). There are three servers built with the official Python SDK, each with Pydantic input
schemas, declared output schemas and \texttt{readOnly}/\texttt{destructive}/\texttt{idempotent} annotations:
\emph{Portfolio} (read-only: \texttt{search\_evidence}, \texttt{get\_case\_study}, \texttt{list\_skills}, resources
\texttt{portfolio://case/\{id\}}, \texttt{portfolio://proposal/\{id\}}, three prompts); \emph{Marketplace} (8 read and 3 write tools,
a required idempotency key, \texttt{dry\_run}, a per-user JWT bearer on Streamable HTTP, a polite per-user limiter of 50/min
against GigSim's 120/min, and an append-only audit log); and \emph{Workspace} (free slots, tentative holds, email drafts on a
documented SQLite mock). Third-party servers are the official Upwork MCP server (OAuth~2.1 with PKCE, tokens Fernet-encrypted
per user) and the GitHub MCP server for repository ingestion.

\textbf{Marketplace port (R1).} \file{agent/marketplace_port.py} defines the interface. \file{mcp_servers/marketplace/adapters/}
holds \texttt{GigSimAdapter} (REST, with bounded retries and backoff that honours \texttt{Retry-After}) and
\texttt{UpworkLiveAdapter} (\texttt{MultiServerMCPClient}; its tool-name mapping lives in config and is validated against
\texttt{list\_tools}). Switching \texttt{marketplace.mode} changes only which adapter the server instantiates. The graph sees
identical tools.

\textbf{Durability and failure handling.} Postgres holds the checkpointer and the store. Approvals resume the interrupted
thread from the UI, and the pending payload is read from the checkpoint, so a page refresh never loses it. Submit
idempotency keys are a pure function of (user, job, draft version). Every node runs under a guard
(\file{agent/graphs/_guard.py}) that applies step/token budgets (``needs attention''), bounded retries with backoff and
timeouts, \emph{degraded mode} (\texttt{MarketplaceUnavailable} moves the job to \texttt{queued\_degraded}, the UI shows a banner
and the worker retries), and a dead-letter status. Background work runs on Arq/Redis with deterministic job ids. A
reconciler re-enqueues runs that were in flight when a worker died.

\section{R\&D question A -- tool granularity}
\textbf{Hypothesis.} Fewer, coarser tools reduce tool-selection errors for a 3B model, while rewritten descriptions recover
argument accuracy. \textbf{Setup} (\file{evals/runners/tool_granularity.py}): 40 tasks with known correct calls, run against
three tool sets (fine-grained, coarse, and best-with-rewritten-descriptions), measuring tool-selection accuracy, argument
accuracy and tokens per task.
\textbf{Result:} \pending. This experiment needs about 120 local LLM tool-calling turns, and it had not completed on the
CPU at this commit. \file{evals/results/tool_granularity.json} will hold the measurement; the conclusion is deliberately
left unwritten until it exists.

\section{R\&D question B -- retrieval, claim verification and what ships}\label{sec:rag}
\textbf{Hypothesis (stated before running).} Requirement decomposition plus hybrid retrieval raises evidence Recall@5 by
$\ge$10 points over V0 on jobs with four or more required skills, because implied skills live in code chunks that whole-post
dense retrieval misses.

\textbf{Setup.} Every variant is a preset of one composable pipeline (\file{rag/pipeline.py}), so the ablation removes one
flag at a time. Embeddings use bge-small-en-v1.5 and the cross-encoder is ms-marco-MiniLM-L6 (both ONNX on CPU). The gold
relevance of a project is the overlap between the job's skills and the project's explicit and implied skills, with
ubiquitous skills (Python, SQL, \ldots) weighted 0.25. Metrics are Recall@5, MRR and nDCG@10 over projects.

\InputIfFileExists{generated/rag_table.tex}{}{\pendingtable{RAG variant comparison}{evals/results/rag_variants.json}}
\InputIfFileExists{generated/rag_breakdown.tex}{}{}
\InputIfFileExists{generated/ablation_table.tex}{}{}

\textbf{Early negative result.} On a 20-job-per-persona smoke run, plain hybrid retrieval (V2) was \emph{worse} than naive
dense retrieval (Recall@5 0.533 vs.\ 0.558). BM25 over code favours files whose identifiers are common words. This is why
the shipped pipeline adds reranking and multi-vector aggregation rather than relying on hybrid retrieval alone.

\textbf{Claim verification (B3)} (\file{rag/verify.py}). After drafting, each first-person experience or result sentence is a
claim. It is labelled supported, partly supported or unsupported using the cross-encoder score against the cited (else
retrieved) evidence, plus entity checks: every skill and every number in the claim must appear in the evidence.
Unsupported claims are removed before the Critic sees the draft. The unsupported-claim \emph{rate} is measured by a
separate, stricter checker built on the persona's ground-truth facts (\file{evals/checkers/fact_check.py}), so the
system is not graded by its own verifier. Before/after rates per proposal: \pending{} (\file{evals/results/proposals.json}).

\textbf{Shipped choice.} The deployed variant is set in \texttt{rag.variant} and must be the best Recall@5/nDCG@10 configuration
in Table~\ref{tab:rag} whose p95 latency fits the analysis budget. The \texttt{search\_evidence} tool keeps the shipped flags but
disables query decomposition, because the Fit Matcher already decomposes at the agent level and doing it twice cost
16--25\,s per call under load.

\textbf{Proposal-level results} (judge score, GigSim reply rate, p95 latency and cost per proposal for V0 vs.\ shipped):
\pending{}. The LLM judge must be calibrated on 30 items labelled by the author (Cohen's $\kappa$); the labelling sheet is
\file{evals/datasets/judge_labels_template.csv}. No judge number is reported without that calibration, and the blind
pairwise classmate panel was not run.

\section{R\&D question C -- is the multi-agent design worth it?}
\textbf{Setup} (\file{evals/runners/arch_compare.py}): (a) one ReAct agent holding all tools
(\file{agent/prompts/react_agent.md}), (b) a \texttt{langgraph-supervisor} with the same workers, and (c) the hand-built
graph, all compared on the same jobs. \textbf{Hand-built graph decisions on the full dev set:}
\InputIfFileExists{generated/decision_table.tex}{}{\pendingtable{Decision quality}{evals/results/decision_eval.json}}
\noindent The comparison with (a) and (b) is \pending: each ReAct job takes 5--8 local LLM turns ($\approx$5\,min on this CPU). The
expected trade-off, to be confirmed or refuted by the measurement, is that the graph spends no LLM tokens on the 70\% of
jobs that should be skipped, while a single agent can over-claim because nothing sits between drafting and output.

\section{The change request}
GigSim's 2 November policy update was not available to me, so I \emph{simulated} one and label it as such (\file{gigsim/app.py},
policy v2). It has two parts. (1) A new rule: proposals may not contain external links (submission returns HTTP 422).
(2) A new cost: premium jobs (Expert level or a budget $\ge$\$1{,}000) cost 4 extra Connects. Absorbing it required no code
change in the agents. The Compliance Guard reads \texttt{policy.v2\_rules} from config (already present, enabled by
\texttt{marketplace.policy\_version: 2}), and Pricing reads the Connects cost that the marketplace returns per job. The diff
is one line of \file{agent/config.yaml} plus the policy-pack text (see \file{docs/CHANGE_REQUEST.md}). This holds because
rules and prompts were in configuration from the start; the risk is that a change requiring new \emph{logic} (e.g.\ a
per-week proposal quota) would need code in \texttt{decide}.

\section{Signature extension -- X2 proposal experiments}
Two proposal strategies (problem-first vs.\ question-first openings) run as an A/B test in GigSim, with a two-proportion
$z$-test and Fisher's exact test on reply rate (\file{evals/runners/ab_test.py}, \file{docs/EXTENSION.md}). The opening style
of every submitted proposal is recorded in the pipeline table and feeds outcome memory, so a significant winner becomes a
promoted insight that the Writer reads. Results: \pending.

\section{Failure analysis}
The taxonomy (missed evidence, wrong skip, mispriced, unsupported claim, generic opening, compliance miss, injection, tool
error) is applied to the rows of \file{evals/results/decision_eval_rows.csv} and the proposal runs. Failures observed during
development, with the fix and its outcome:
\begin{enumerate}
\item \textbf{Wrong skip / borderline fit.} A Davinci proof-of-concept job scored fit 0.40 because sentence-only requirements
  (``build a functional model\ldots'') counted as much as skills. Fix: such requirements weigh 0.5. Re-measured in the decision eval.
\item \textbf{Gold-label error.} A Solidity/NFT job was labelled ``bid'' for the full-stack persona because React and JS covered 2 of 3
  skills. Fix: title skills are core and must all be covered. Bid share fell from 32--40\% to 25--30\%.
\item \textbf{Tool error (degraded mode).} Port collision: another local service answered GigSim's port with HTML, which the
  adapter tried to parse as JSON. Fix: tolerate non-JSON errors and raise a typed \texttt{MarketplaceError}.
\item \textbf{Tool error (lifecycle).} MCP sessions entered in one task and closed in another crashed anyio. Fix: one owner task per session.
\item \textbf{Missed evidence.} Cross-encoder relevance between ``Celery workers with Redis'' and code chunks was $\approx$0. Fix:
  requirement strength also counts skill metadata extracted from code imports.
\item \textbf{Latency.} \texttt{search\_evidence} took 16--25\,s because the shipped variant re-decomposed atomic requirements. Fix: a
  single-requirement retrieval path.
\item \textbf{Mis-detection risk.} Injection patterns also match benign text such as ``new instruction manual''. The false-positive
  rate is measured in Table~\ref{tab:dec} (injection column) rather than assumed.
\end{enumerate}
Further failures from the proposal runs (unsupported claims, generic openings, compliance misses): \pending.

\section{Safety and platform respect}
\textbf{Never does:} submit, message or spend Connects without an approval interrupt; scrape, automate a browser or use session
cookies; claim experience the evidence does not support; follow instructions found inside a job post or client message.
\textbf{Enforcement:} untrusted text is PII-masked and delimited at intake, and spoofed delimiters are neutralised. The Risk
Assessor detects injected instructions and those sentences are stripped before analysis. The Writer prompt forbids
following them, and Compliance removes any echo (e.g.\ ``BANANA'', ``\$5'', guarantees). Writes require approval plus allowlists
plus idempotency keys. PII masking covers logs, the audit log and LangSmith hide hooks, and OAuth tokens are encrypted and
never logged. Every query is tenant-scoped by \texttt{user\_id}, and the audit table is append-only via a trigger.
\textbf{What still gets through:} paraphrased injections that match no pattern reach the Writer as delimited data; only the
prompt and the compliance echo check defend there, and the attack success rate before and after defences is \pending. The
cross-encoder is a relevance model, not an entailment model, so a fluent but subtly wrong claim near true evidence can pass
as ``partly supported''.

\section{Hidden-test results against declared budgets}
The instructor's hidden run cannot be reproduced here. The stand-in is the git-ignored held-out split (80 jobs per persona,
new clients), run with \texttt{--split test}. Results: \pending. Analysis latency on the dev set: p95 \decPninetyfive\,s against a
20\,s budget (Table~\ref{tab:dec}). Proposal latency and cost: \pending. The \textbf{deployed link} was not live at this commit:
Render/Railway deployment requires the author's account, and the hosted-open profile requires a Groq API key, which was not
available. The open (local) profile is the one that runs.

\section{Models: open configurations compared}\label{sec:models}
Two open configurations are defined (\file{agent/config.yaml}): \texttt{open} (Qwen2.5-3B for every node, with the 7B as judge)
and \texttt{open\_7b\_writer} (7B Writer). Measured throughput on the i7-8550U: Qwen2.5-3B 6.1 tok/s generation and 13.1 tok/s
prompt (cold); Qwen2.5-7B 2.9 and 8.0 tok/s. The quality comparison between the two configurations is \pending. Embeddings and
reranker are open ONNX models on CPU in every profile.

\appendix
\section{Reproducing every number}
\begin{verbatim}
make data          # Kaggle -> personas, GigSim seed, dev/held-out sets (seeded)
make eval-rag      # evals/results/rag_variants.json (+ _rows.csv)
make eval-decision # evals/results/decision_eval.json
make eval          # all of the above + LLM evals
python -m scripts.build_report_tables && cd report && latexmk -pdf report.tex
\end{verbatim}
\section{AI usage}
See \file{docs/AI_USAGE.md}. An AI coding assistant (Claude Code) helped write most of the code under my direction. One
error it made that was caught: its first data builder applied about 500 separate regexes per job across 110k jobs and ran
for over 10 minutes; it was replaced by a substring-prefiltered matcher (97\,s end to end).
\end{document}
